\tips

\section*{Tips for the Reader}
Every chapter begins with a brief introduction, a mini table of contents and objectives. It is recommended that you read these and then skim through the rest of the chapter. Every chapter ends with a chapter summary that includes the \textit{Next Steps} to help students understand what has been covered and where we are headed. Chapters also end with a few sample questions for review. During your pre-read skim you should make mental note of the headings, subheadings, figures, and equations. The purpose of skimming is to get a sense of what is coming up as you start to read the chapter. When you start reading, read as if you are listening to the author speak. If you get a question in your head as your listening to me, write it down. Proceed to read even if you're not completely sure you understand what you're reading. Sometimes a concept becomes more clear as you proceed. If not, make a note of what you're not understanding and bring it to class. Identifying and clearing confusion is one of the purposes of class time, the second purpose of class time is application, depth and breadth.

\section*{A Note about Chapter Objectives}

The chapter objectives include specific behavioral terms. There is an ordered relationship between the behavioral terms that represents increasing and nested expectations. 

To \textbf{define} is to know the words being used, to know what they represent. Whether you can define something can be tested with a straightforward question such as identifying the correct definition from a set of definitions; whether you can write the definition; respond true or false to a proposed definition for a term; or identify the correct term from a set of terms when presented with a definition. 

To \textbf{explain} goes beyond defining, but it also implies that you can define the required terms. Explaining is the act associated with understanding. If you understand something you can explain it. Objectives are typically written as acts - what the person learning can do if they have achieved the objective. Since understanding is observed as an act of explaining, we will use the objective "explain" when there is something you must "understand". How do you demonstrate that you understand? You explain whatever it is that you understand. To explain you simultaneously hold a set of definitions together including their relationship to one another. You can explain a model by demonstrating that you know all the definitions involved, and describe how they relate to one another, and can communicate that explanation. Whether you can explain something can be tested by asking about several aspects of what you are to explain or asking you to infer which part is missing; or by asking about one part and asking you to infer the other part. The point is, to explain goes beyond defining and has an expectation of definition. But to define does not require you to explain. 

To \textbf{evaluate} requires you to make judgements and to infer beyond the thing you are asked to evaluate. To evaluate implies that you can define and explain the thing (or concept) you are asked to evaluate. Being tested for whether you can evaluate can include asking you to infer to or from something that has not been directly covered or discussed, going beyond the thing (or concept) you are asked to evaluate to consider its implications in different situations. 

\textbf{Providing an example} is an objective based on a specific task and is considered a form of asking you to evaluate since it goes beyond a particular thing (or concept) to a totally new thing (or concept). Asking you to provide an example of a model requires you to know the concept of a model well enough that you can use your background knowledge to come up with something you're already familiar with that fits what you know about models. That level of knowledge requires you to evaluate a model because you evaluated your background knowledge and determined that it is a model. If a chapter objective asks you to provide an example of a model that is useful to physical therapy practice it assumes you will be able to define, explain and evaluate the concept of a model, and that you have background knowledge about physical therapy practice that you can use to evaluate models used in practice to consider what those models are and whether they are useful. Providing an example is similar to the objective \textbf{Apply}. With both of these objectives you are being asked to use your knowledge of the general concept (or thing) to a specific instance. If the concept was fruit, then an example is an apple. And you can apply the the concept of fruit by listing several fruits including what makes them all fruits.

\textbf{Compare and contrast} is also an objective based on a specific task that requires evaluation of two (or more) things (or concepts) to identify what is alike and different about them. This can occur with the general concept (Fruit) by comparing fruits with vegetables, or with specifics (apples) by comparing different specific fruits (apples and oranges), or by comparing different specifics from each concept (apples and broccoli). Comparing and contrasting is an important higher level evaluative process that helps you to develop a deeper understanding.

\vspace{\baselineskip}
\begin{flushright}\noindent
Boscawen, New Hampshire\hfill {\it Sean  Collins}\\
May 2024\hfill {\it \hspace{1cm} }\\
\end{flushright}


